\documentclass[11pt]{amsart}

\usepackage{amsmath}
\usepackage{amssymb}
\usepackage{amsthm}
\usepackage{url}
\usepackage{hyperref}
%\usepackage{psfig}

\begin{document}
\baselineskip=16pt
\textheight=8.5in
%\parindent=0pt 
\def\sk {\hskip .5cm}
\def\skv {\vskip .05cm}
\def\cos {\mbox{cos}}
\def\sin {\mbox{sin}}
\def\tan {\mbox{tan}}
\def\intl{\int\limits}
\def\lm{\lim\limits}
\newcommand{\frc}{\displaystyle\frac}
\def\xbf{{\mathbf x}}
\def\fbf{{\mathbf f}}
\def\gbf{{\mathbf g}}

\def\dbA{{\mathbb A}}
\def\dbB{{\mathbb B}}
\def\dbC{{\mathbb C}}
\def\dbD{{\mathbb D}}
\def\dbE{{\mathbb E}}
\def\dbF{{\mathbb F}}
\def\dbG{{\mathbb G}}
\def\dbH{{\mathbb H}}
\def\dbI{{\mathbb I}}
\def\dbJ{{\mathbb J}}
\def\dbK{{\mathbb K}}
\def\dbL{{\mathbb L}}
\def\dbM{{\mathbb M}}
\def\dbN{{\mathbb N}}
\def\dbO{{\mathbb O}}
\def\dbP{{\mathbb P}}
\def\dbQ{{\mathbb Q}}
\def\dbR{{\mathbb R}}
\def\dbS{{\mathbb S}}
\def\dbT{{\mathbb T}}
\def\dbU{{\mathbb U}}
\def\dbV{{\mathbb V}}
\def\dbW{{\mathbb W}}
\def\dbX{{\mathbb X}}
\def\dbY{{\mathbb Y}}
\def\dbZ{{\mathbb Z}}


%Calligraphic roman letters
\def\calA{{\mathcal A}}
\def\calB{{\mathcal B}}
\def\calC{{\mathcal C}}
\def\calD{{\mathcal D}}
\def\calE{{\mathcal E}}
\def\calF{{\mathcal F}}
\def\calG{{\mathcal G}}
\def\calH{{\mathcal H}}
\def\calI{{\mathcal I}}
\def\calJ{{\mathcal J}}
\def\calK{{\mathcal K}}
\def\calL{{\mathcal L}}
\def\calM{{\mathcal M}}
\def\calN{{\mathcal N}}
\def\calO{{\mathcal O}}
\def\calP{{\mathcal P}}
\def\calQ{{\mathcal Q}}
\def\calR{{\mathcal R}}
\def\calS{{\mathcal S}}
\def\calT{{\mathcal T}}
\def\calU{{\mathcal U}}
\def\calV{{\mathcal V}}
\def\calW{{\mathcal W}}
\def\calX{{\mathcal X}}
\def\calY{{\mathcal Y}}
\def\calZ{{\mathcal Z}}


\def\la{{\langle}}
\def\ra{{\rangle}}
\def\summ{{\sum\limits}}
\def\eps{{\varepsilon}}
\def\lam{{\lambda}}
\def\tr{{\rm tr}}
\def\char{{\rm char}}
\def\Mat{{\rm Mat}}
\def\Func{{\rm Func}}
\def\Span{{\rm Span}}
\def\Spec{{\rm Spec}}
\def\supp{{\rm supp}}
\def\Inv{{\rm Inv}}
\def\inv{{\rm inv}}
\def\Ker{{\rm Ker}}
\def\Im{{\rm Im}}
\def\rk{{\rm rk}}
\def\phi{{\varphi}}
\def\row{{\rm row}}
\def\col{{\rm col}}
\def\Rowspace{{\rm Rowspace}}
\def\null{{\rm null}}


\bf\centerline{Homework \#10. Due on Friday, Nov 20th, 11:59pm on Canvas}\rm
\vskip .1cm

\bf{Reading for this homework: }\rm parts of 5.4, 7.1 and 7.3

\bf{Plan for next week (Nov 16,18): }\rm  Jordan Canonical Form (continue 7.1 and 7.2)). 

\skv
\noindent 

\bf\centerline{Problems: }\rm
\skv
{\bf Problem 1:} Let $F$ be a field and $A_1,\ldots, A_k$ square
matrices over $F$ (possibly of different sizes). Suppose that  $A_i\in \Mat_{n_i}(F)$, and
let $n=\sum\limits_{i=1}^k n_i$. Denote by $A_1\oplus\ldots\oplus A_k$ the 
block-diagonal matrix whose diagonal blocks are $A_1,\ldots, A_k$, in this order (thus $A_1\oplus\ldots\oplus A_k\in \Mat_n(F)$).
Prove that if $A=A_1\oplus\ldots\oplus A_k$, then 
\begin{itemize}
\item[(a)] $\rk(A)=\sum\limits_{i=1}^k \rk(A_i)$;
\item[(b)] for any polynomial
$p(x)\in\mathcal P(F)$ we have $p(A)=p(A_1)\oplus\ldots\oplus p(A_k)$.
\end{itemize}
See a hint for (b) at the end of the homework.
\skv

{\bf Problem 2:} Let $V$ be a finite-dimensional vector space, $T\in\calL(V)$ and $W$ a $T$-invariant subspace of $V$.
Let $T_W\in \calL(W)$ be the restriction of $T$ to $W$.
\begin{itemize}
\item[(a)] Prove that $\chi_{T_W}(x)\mid \chi_{T}(x)$ and $\mu_{T_W}(x)\mid \mu_T(x)$ (Recall that $\mu_S(x)$ is the minimal polynomial
of a linear map $S$).
\item[(b)] Use (a) (and something we proved in class) to show that if $T$ is diagonalizable, then $T_{W}$ is also diagonalizable.
\end{itemize}
\skv
{\bf Problem 3:} Let $V$ be a finite-dimensional vector space over a field $F$ and $T\in\calL(V)$. Recall that for each $\lam\in \Spec(T)$,
the generalized $\lam$-eigenspace is $K_{\lam}(T)=\bigcup_{k=0}^{\infty}\Ker((T-\lam I)^k)$. In class we proved that
\begin{itemize}
\item[(i)] the sum $\sum\limits_{\lam\in\Spec(T)}K_{\lam}(T)$ is always direct (Theorem 21.1);
\item[(ii)] if $\chi_T(x)$ splits over $F$, then $\sum\limits_{\lam\in\Spec(T)}K_{\lam}(T)=V$ (Theorem 21.3; proof will be completed in
Lecture~22).
\end{itemize}
The goal of this problem is to provide an alternative proof of (i) assuming (ii). This proof only works under the assumption that 
$\chi_T(x)$ splits, but once (i) is proved under that assumption, it is not hard to deduce (i) in general using the argument we used
in the proof of the Cayley-Hamilton Theorem in Lecture~20.
\skv
So assume that $\chi_T(x)$ splits over $F$.
\begin{itemize}
\item[(a)] Given $\lam\in \Spec(T)$, let $T_{\lam}\in \calL(K_{\lam})$ be the restriction of $T$ to $K_{\lam}$. Prove that
$\chi_{T_{\lam}}(x)=(x-\lam)^{\dim(K_{\lam})}$ (this can be done in several different ways).
\item[(b)] Use (a) and Problem~2 to prove that for each $\lam\in\Spec(T)$ we have $\dim K_{\lam}(T)\leq m_{\lam}(T)$
where $m_{\lam}(T)$ is the algebraic multiplicity of $\lam$.
\item[(c)] Now use (b), (ii) above and HW\#8.6 to prove that $\dim K_{\lam}(T) = m_{\lam}(T)$ for each $\lam$ and the sum $\sum\limits_{\lam\in\Spec(T)}K_{\lam}(T)$ is direct.
 \end{itemize}
\skv
{\bf Problem 4:} This problem strengthens the statement of 
Problem~4(b) in HW\#9. Let $V$ be a vector space and
$T\in \mathcal L (V)$.
\begin{itemize}
\item[(a)] Suppose that there exists $i\in\dbN$ be such that $\Ker(T^{i+1})=\Ker(T^{i})$. Prove that
$\Ker(T^{m})=\Ker(T^{i})$ for all $m>i$.  Note that we are not assuming that $V$ is finite-dimensional.
{\bf Hint:} It is enough to show
that $\Ker(T^{i+2})=\Ker(T^{i+1})$.
\end{itemize}
In parts (b) and (c)  assume that $V$ is finite-dimensional, and let $n=\dim V$.

\begin{itemize}
\item[(b)] Prove that $\cup_{i=0}^{\infty}\Ker(T^{i})=\Ker(T^n)$.
\item[(c)] Deduce from (b) that $K_{\lam}(T)=\Ker((T-\lam I)^n)$ for any $\lam\in \Spec(T)$.
\end{itemize}
\skv
{\bf Problem 5:} Let $V$ be a vector space over a field $F$ and $T\in\calL(V)$.
\begin{itemize}
\item[(a)] Prove that if $W_1,\ldots, W_k$ are $T$-invariant subspaces of $V$, then $\sum_{i=1}^k W_k$ is also $T$-invariant.
\item[(b)] Assume that $\Spec(T)$ is finite (of course, this is automatic if $V$ is finite-dimensional), and let $m=|\Spec(T)|$.
Use (a) to prove that $V$ has at least $2^m$ distinct $T$-invariant subspaces.
\item[(c)] Now assume that $n=\dim(V)$ is finite and that $T$ has $n$ distinct eigenvalues $\lam_1,\ldots, \lam_n$.
This means that $T$ is diagonalizable and each eigenspace $E_{\lam_i}(T)$ is $1$-dimensional. Choose a nonzero $w_i\in E_{\lam_i}(T)$,
for each $i$, and let $\beta=\{w_1,\ldots,w_n\}$ (thus, $\beta$ is a basis consisting of eigenvectors for $T$).

Given $v\in V$, write $v=\sum_{i=1}^n \lam_i w_i$, and define $$\supp(v)=\{1\leq i\leq n: \lam_i\neq 0\}.$$ In other words,
$\supp(v)$ is the set of all $i$ such that the $i^{\rm th}$ coordinate of $[w]_{\beta}$ is nonzero.
Prove that the smallest $T$-invariant subspace containing $v$ is equal
to $\sum_{\lam\in \supp(v)}E_{\lam}(T)$. {\bf Hint:} one way to prove this is using a Vandermonde-type argument similar
to the second proof of the theorem that the sum of eigenspaces is direct (Lecture~18).
\item[(d)] Let $T$ be as in (c). Use (c) and your solution to (b) to prove that $T$ has precisely $2^n$ $T$-invariant subspaces.
\end{itemize}
\skv
{\bf Problem 6:} Let $V$ be a vector space and $T\in\calL(V)$. A sequence of vectors $v_1,\ldots, v_k\in V$
is called a \emph{nilpotent $T$-cycle} if each $v_i$ is nonzero, $T(v_1)=0$ and $T(v_i)=v_{i-1}$ for $2\leq i\leq k$.
\begin{itemize}
\item[(a)] Prove that the vectors which from a nilpotent $T$-cycle must be linearly independent.
\item[(b)] Let $v_1,\ldots, v_k$ be a nilpotent $T$-cycle, $\beta=\{v_1,\ldots, v_k\}$ and $W=\Span(\beta)$ (so that
$\beta$ is a basis of $W$ by (a)). For each $1\leq i\leq k$ let $W_i=\Span(v_1,\ldots, v_i)$; also set $W_0=0$.
Prove that each $W_i$ is $T$-invariant and each $T$-invariant subspace of $W$ is equal to $W_i$ for some $i$.  
\end{itemize}

\newpage 
{\bf Hint for 1(b):} Let $v\in F^n$ and write $v=\sum_{i=1}^k v_i$ where $v_i$ agrees with $v$ in the coordinates corresponding to the
$i^{\rm th}$ blocks and all other coordinates of $v_i$ are equal to $0$. Start by expressing $Av$ in terms of $A_i$ and $v_i$.

\end{document}
